\section{Introduction}\label{sec:introduction}
Encryption conceals message content, yet LLM engines and agents expose packet sizes, cache timing, visited domains and tool sequences \cite{weiss2024prompt,promptpeek2025,earlybird2025,agentprint2025,uitraces2026,jeong2026wra}. A defense result applies to a particular receiver, execution boundary and workload. Reusing a headline score across a frontend, dispatcher or tool provider can change both its protected signal and its cost object.

Our SoK uses the \emph{evaluated result record} to make these transfers inspectable. Separating Dangol's stored-capture and live-protocol results prevents subtracting unmatched scores; separating Early Bird's KV-prefix and semantic-cache mechanisms prevents pricing one with the other's overhead. Combining signal origin with observer access and threat polarity explains why read-only privacy cover and state-changing bot evasion have different utility consequences. These findings emerge from binding evidence across records that a paper-level summary combines.

\paragraph{Central question and contributions}
\emph{Which changes invalidate the transfer of a measured defense result?} The corpus map locates the systems being compared; result records bind each claim to an operating point; the audit identifies the fields that must be checked when that point changes. These are components of one transfer analysis.
\begin{itemize}
\item \textbf{An evidence contract for transfer.} Source-bound records separate receiver, state, target, attacker, utility and cost, making unsupported transfers visible within and between study families.
\item \textbf{Boundary-breaking demonstrations.} Exp~1 changes the emitter; Exp~2 moves cover across the read/write state boundary; Exp~3 evaluates training-calibrated defenses on held-out customers; Exp~4 changes observation and environment from length-only/offline comparisons to measured timing and emulated WAN conditions.
\item \textbf{Operational decisions.} Re-measure the final emitter, isolate state-changing cover, freeze calibration before held-out evaluation, and refit a timing-aware attacker before carrying a protection claim into a new environment. Two-endpoint attacker and calibration matrices show why transfer must be checked separately for each target and direction.
\end{itemize}
\begin{table*}[t]\centering\footnotesize
\caption{Research questions, systematization findings and evidence. Experimental follow-ups test the comparison rules; the SoK contribution is the record-level synthesis.}\label{tab:rq-findings}
\begin{tabularx}{\textwidth}{@{}p{.20\textwidth}YY@{}}\toprule
Question & Finding & Evidence in the main paper\\\midrule
RQ1: targets and observables & A common signal can serve confidentiality, attribution, evasion or authorized diagnostics; polarity changes the protected interest. & Complete core-family map (\S\ref{sec:corpus-map}, Table~\ref{tab:core-family-map}).\\
RQ2: origin versus access & B-size cover does not bind A-call order; plaintext re-emission creates a fresh B trace; shared A-call projections do not establish unchanged state. These distinctions identify the evidence required to reuse a result. & NetEcho/SplitDecoy, WRA/AHB and dispatcher read/write records (\S\ref{sec:discussion}); Exps~1--2.\\
RQ3: portability & A task holdout does not establish topology, hardware or model transfer; mixed configurations can preserve shortcuts. & Generator-specific holdouts (\S\ref{sec:evaluation}); customer-grouped analysis (\S\ref{sec:r5}).\\
RQ4: protection, state and cost & Bind inference, utility and cost to one scoped operating point. The 20-family audit gives R8 Bound 3/17 for component-or-task utility and 2/16 under a task-only sensitivity. & Twenty-family R1--R9 matrix; Exps~3--4 leakage/state/cost contrasts (\S\ref{sec:controlled-placement}). The same audit locates the utility boundary of our demonstrations.\\\bottomrule
\end{tabularx}\end{table*}


\soktakeaway{Scope of claims}{We contribute record-level systematization and controlled demonstrations of evidence transfer. The real-agent target is a requested operation category; the known retail cohort contains 53 customers. Measurements establish fitted-suite inference, logical state conformance and scoped resource costs. A separate 100-episode live study measures prefix-shadow utility under injected deadlines. Deployment-wide security and full joint utility require the validation specified in \S\ref{sec:validity}.}

Figure~\ref{fig:system-model} places trace and action paths on the same agent--tool chain. This joint view makes the receiving representation and the execution boundary explicit before selecting a defense.
\begin{figure*}[t]
\centering\begin{tikzpicture}[font=\scriptsize,>=Latex,
box/.style={draw,rounded corners,align=center,text width=3.35cm,minimum height=10mm,fill=blue!5},
boundary/.style={box,fill=orange!12},
observer/.style={box,fill=green!8}]
\node[box] (llm) at (0,0) {LLM service\\token stream / cache activity};
\node[boundary] (api) at (4.25,0) {Upstream protected link\\padding or batching};
\node[boundary] (front) at (8.5,0) {Content-aware frontend\\decoded content $\rightarrow$ new stream};
\node[observer] (net) at (12.75,0) {Final network observer\\groups, lengths, timing};
\draw[->] (llm)--(api);
\draw[->] (api)--(front);
\draw[->] (front)--(net);
\node[box] (policy) at (0,-2.1) {Agent action policy\\task goal / current state};
\node[boundary] (cover) at (4.25,-2.1) {Cover boundary\\discard, shadow, or execute};
\node[box] (dispatch) at (8.5,-2.1) {Tool / GUI dispatch\\real and admitted cover actions};
\node[observer] (state) at (12.75,-2.1) {Environment / receiving site\\events and changed task state};
\draw[->] (policy)--(cover);
\draw[->] (cover)--(dispatch);
\draw[->] (dispatch)--(state);
\draw[->,dashed] (llm.south)--node[left,align=right]{action\\selection}(policy.north);
\node[font=\scriptsize,align=center] at (6.37,-.95)
{Trace path: upstream protection must be related to the final emitting interface (NetEcho).};
\node[font=\scriptsize,align=center] at (6.37,-3.0)
{Action path: executed cover can change state (AHB); timestamp-only edits can retain event structure (UI Traces).};
\end{tikzpicture}

\caption{Observation and execution boundaries. The four demonstrations test emitter placement, stateful execution/link composition, real-agent operation signatures, and timing with matched transport controls.}
\label{fig:system-model}
\end{figure*}
